\begin{lemma}[Height-five gamma-edge configuration]
Let $N$ be normal globally maximum data for a $k$-triple
$(T,e_0,\varphi)$, suppose that the chosen height of $N$ and the Scheide
height $h(G)$ are both five, and write the five trunk vertices as
$y_i=T.\mathsf{vertex}(i)$ for $0\leq i\leq4$.  Let
$\gamma\in\overline\varphi(y_0)$.  Suppose that specified identities
$g\in E(G)$ and $v\in V(T)\setminus\{y_0,y_1,y_2,y_3,y_4\}$ satisfy
$\varphi(g)=\gamma$ and
$\operatorname{ends}(g)=\{T.\mathsf{vertex}(i),v\}$ for at least one
$i\in\{0,1,2,3,4\}$.  Then there are edge identities $f_1,f_2,f_3$ and
vertices $v_1,v_2\in V(T)\setminus\{y_0,y_1,y_2,y_3,y_4\}$ such that
\[
 \operatorname{ends}(f_1)=\{y_1,v_1\},\qquad
 \operatorname{ends}(f_2)=\{y_4,v_2\},\qquad
 \operatorname{ends}(f_3)=\{y_2,y_3\},
\]
and
$\varphi(f_1)=\varphi(f_2)=\varphi(f_3)=\gamma$.
The three displayed objects are edge identities; parallel edges are not
identified with their unordered endpoint pairs.
\end{lemma}
\begin{proof}
Put $y_i=T.\mathsf{vertex}(i)$ for $0\leq i\leq4$ and
$q=N.\mathsf{closingEdge}$.  The equality
$N.\mathsf{height}=5$, the height upper bound, the closing-index equation,
and the support and endpoint fields in
\noderef{NormalGlobalMaximumTriple} show that the stored closing path is
\[
 y_1,e_2,y_2,e_3,y_3,e_4,y_4,q,y_0.                 \tag{1}
\]
Here the paper's $e_j$ is the Lean identity $T.\mathsf{edge}(j-1)$, and
\noderef{TreeSequence} gives $e_1=e_0$ and makes the five displayed
vertices distinct.  In particular, (1) uses the individual identities
$e_2,e_3,e_4,q$, even if another edge has the same endpoint pair.

The global field of $N$ is the condition
\noderef{GlobalMaximumTriple}.  Its \noderef{KTriple} clause says that
$\varphi$ is proper, leaves exactly $e_0$ uncolored, and has all the stated
criticality, index, and degree data.  Each edge-node in (1) belongs to the
incidence graph \noderef{TwoColorIncidenceGraph} for the colors
$\alpha=N.\mathsf{alpha}$ and $\beta=N.\mathsf{beta}$, so each is colored
$\alpha$ or $\beta$.  Since $\alpha$ is missing at $y_0$, the edge $q$
incident with $y_0$ has color $\beta$.  Successive applications of
properness, in the sense of \noderef{IsProperEdgeColoring}, along (1) then
give
\[
 \varphi(e_2)=\varphi(e_4)=\alpha,
 \qquad \varphi(e_3)=\varphi(q)=\beta.              \tag{2}
\]

Apply \noderef{ScheideKempePathHeightBound} to the globally maximum triple,
the color $\gamma$ missing at $y_0$, and the color $\beta$ missing at
$y_1$.  It gives a simple incidence path $P_1$ from $y_0$ to $y_1$ whose
support is the full component
\noderef{TwoColorComponent} for $\gamma,\beta$, all of whose multigraph
vertices lie in $V(T)$.  It also gives
\[
 \operatorname{length}(P_1)/2+1\leq h(G)=5.          \tag{3}
\]
The length here is incidence-walk length.  Because the path begins and ends
at vertex-nodes and alternates between vertex-nodes and edge-nodes, the
left side of (3), not the raw incidence length, is the number of multigraph
vertices of $P_1$.

We first justify that the two colors defining this path are distinct.  The
vertex-injectivity clause of \noderef{TreeSequence} gives $y_0\ne y_1$.
Thus the $y_0$--$y_1$ path $P_1$ has positive length, so it has a first
edge-node, say $d$, adjacent to its initial vertex-node $y_0$.  Expanding
this adjacency with \noderef{TwoColorIncidenceGraph} gives
$\noderef{Incident}(G,y_0,d)$ and says that
$\varphi(d)=\gamma$ or $\varphi(d)=\beta$.  If $\gamma=\beta$, either
color alternative says $\varphi(d)=\gamma$.  The incidence then puts
$\gamma$ in the set of colors present at $y_0$ by
\noderef{PresentColors}.  This contradicts
$\gamma\in\overline\varphi(y_0)$, because \noderef{MissingColors} defines
the latter set as the complement of the former.  Hence
$\gamma\ne\beta$.

The identity $q$ has color $\beta$ and is incident with $y_0$, so its
edge-node, and then its other endpoint $y_4$, lie in the full component
$P_1$.  Hence $P_1$ contains the three distinct multigraph vertices
$y_0,y_1,y_4$.  We also record that
\[
 \gamma\ne\alpha.                                   \tag{4}
\]
Indeed, if the two colors were equal, the normal path (1) would lie in the
same full $\gamma,\beta$ component as $P_1$.  It already contains all five
trunk vertices, while the assumed $\gamma$-edge $g$ connects one of those
vertices to the distinct outside vertex $v$.  The component would therefore
contain at least six multigraph vertices, contradicting (3).  The
individual identity $g$ remains visible as an edge-node here; the argument
uses its stated endpoint equality and does not identify it with any other
edge having the same ends.

We next prove that no $\gamma$-edge joins either $y_2$ or $y_3$ to a
vertex outside the five-vertex trunk.  Suppose, to the contrary, that an
identity $r$ of color $\gamma$ joins one of $y_2,y_3$ to such a vertex
$w$.  If any one of $y_2,y_3,w$ belonged to $P_1$, then the
$\beta$-identity $e_3$ between $y_2,y_3$ and the $\gamma$-identity $r$
would put all three in the connected component.  Together with
$y_0,y_1,y_4$, these are six distinct multigraph vertices, contrary to
(3).  Thus
\[
 y_2,y_3,w\notin V(P_1).                             \tag{5}
\]

For completeness, we now perform the switch which rules out (5).  Global
maximality makes $T$ fixed-pair maximal in the sense of
\noderef{FixedPairMaximalTashkinovTree}: any fixed-pair Tashkinov tree
containing $V(T)$ has order at most the global maximum
\noderef{TashkinovOrder}, and hence has the same finite vertex set.  The
condition that $\gamma$ is unused through $y_1$ in
\noderef{TreeColorUnusedThrough} is vacuous, since there is no colored tree
edge strictly between indices zero and one.  Therefore
\noderef{ScheideMaximalKempePath}, applied to $y_0,y_1$ and
$\gamma,\beta$, says clause by clause that the component swap
\[
 \varphi'=\noderef{TwoColorComponentSwap}
   (G,\varphi,\gamma,\beta,y_0)
\]
is proper, still has exactly $e_0$ uncolored, and leaves the same ordered
sequence $T$ a Tashkinov tree.  These are precisely the coloring and tree
clauses in \noderef{KTriple}, \noderef{IsOneEdgeUncolored}, and
\noderef{TashkinovTree}; the critical, index, and degree clauses do not
depend on the coloring.  The vertex set and its cardinality have not
changed.  Thus
$(T,e_0,\varphi')$ again satisfies \noderef{GlobalMaximumTriple}.

Here is the required missing-color check after the swap.  The color
$\beta$ was missing at $y_1$.  The last edge-node of $P_1$ is therefore
colored $\gamma$, and properness makes it the only incident
$\gamma$-identity at $y_1$.  The full-component swap changes this identity
to $\beta$ and changes no identity outside the component, so $\gamma$ is
missing at $y_1$ under $\varphi'$.  By (4), $\alpha\ne\gamma$; moreover
$q$ is a $\beta$-edge incident with $y_0$ while $\alpha$ is missing there,
so $\alpha\ne\beta$.  Hence $\alpha$ is not one of the switched colors and
remains missing at $y_0$.  The edge $q$ changes from $\beta$ to $\gamma$,
whereas $e_2,e_4$ retain color $\alpha$ by (2).  Finally, (5) excludes both
endpoints of $r$ from the full component.  If the edge-node $r$ belonged to
it, its incidence adjacency in \noderef{TwoColorIncidenceGraph} would put
an endpoint there as well.  Thus $r$ is not switched and retains color
$\gamma$.

Apply \noderef{ScheideKempePathHeightBound} again, now to the globally
maximum triple $(T,e_0,\varphi')$, with $\alpha$ missing at $y_0$ and
$\gamma$ missing at $y_1$.  Its full $\alpha,\gamma$ incidence path $P_2$
contains the forced segments
\[
 y_0,q,y_4,e_4,y_3,
 \qquad y_1,e_2,y_2.                                 \tag{6}
\]
Since $r$ joins $w$ to one of $y_2,y_3$, (6) also puts the edge-node $r$
and the vertex-node $w$ in $P_2$.  Consequently $P_2$ contains the six
distinct multigraph vertices $y_0,y_1,y_2,y_3,y_4,w$.  The second
height-bound application says that its incidence length $L$ satisfies
$L/2+1\leq5$, again bounding the multigraph-vertex count rather than the
incidence length.  This contradiction proves
\[
 \text{no $\gamma$-edge from $\{y_2,y_3\}$ leaves the trunk.} \tag{7}
\]

It remains to identify the three asserted edge identities.  By
\noderef{ScheideGlobalMaximumOrder}, the finite set
\noderef{TreeSequenceVertexSet} of a globally maximum triple is elementary.
Unfolding this assertion with \noderef{ColoredVertexSetElementary} says
that distinct vertices of this set have disjoint missing-color sets.  Each
$y_i$ belongs to the set by \noderef{TreeSequenceVertexSet}, and the
vertex-injectivity clause of \noderef{TreeSequence} gives $y_i\ne y_0$ for
$1\leq i\leq4$.  Thus $\gamma$ cannot be missing at any of
$y_1,y_2,y_3,y_4$: otherwise it would belong to that vertex's missing set
and to the missing set at $y_0$, contrary to the stated disjointness.  The
definitions \noderef{MissingColors} and \noderef{PresentColors} therefore
supply a $\gamma$-colored incident edge at each of these four vertices.
Properness in \noderef{IsProperEdgeColoring} makes that identity unique at
each vertex.

Let the unique $\gamma$-identities at $y_2,y_3$ be $a_2,a_3$.  For
$i=2,3$, their defining incidence and \noderef{Incident} let us choose
the other endpoint $z_i$ with
$\operatorname{ends}(a_i)=\{y_i,z_i\}$; this equality records both
incidence memberships.  Looplessness gives $z_i\ne y_i$.  By (7), these
other endpoints lie in the trunk, and neither can be $y_0$, where
$\gamma$ is missing.  If $a_2\ne a_3$, the matching property furnished by
properness forces $z_2,z_3$ to be $y_1$ and $y_4$ in
some order: using $y_2$ or $y_3$ would make the two identities coincide,
and two distinct identities cannot share an endpoint by
\noderef{AdjacentEdges}.  The unique $\gamma$-edges at all four vertices
$y_1,y_2,y_3,y_4$ would then be internal to the trunk.  Together with (7)
and the missing color at $y_0$, this leaves no trunk vertex incident with
the assumed crossing identity $g$, a contradiction.  Hence
$a_2=a_3$; call this identity $f_3$.  It is incident with both $y_2$ and
$y_3$.  Expanding the first membership with \noderef{Incident}, write
$\operatorname{ends}(f_3)=\{y_2,z\}$.  Looplessness in
\noderef{LooplessMultigraph} gives $z\ne y_2$.  The second incidence says
$y_3\in\{y_2,z\}$, while vertex injectivity gives $y_3\ne y_2$; hence
$z=y_3$.  Therefore
$\operatorname{ends}(f_3)=\{y_2,y_3\}$ and
$\varphi(f_3)=\gamma$.

Call the unique $\gamma$-edge at $y_1$ by the name $f_1$.  Its incidence
at $y_1$, unfolded by \noderef{Incident}, lets us choose the other
endpoint $v_1$.  Thus both $y_1$ and $v_1$ are incident with this same
identity and
$\operatorname{ends}(f_1)=\{y_1,v_1\}$; looplessness gives
$v_1\ne y_1$.  This other endpoint is not $y_0$, where
$\gamma$ is missing, and is not $y_2$ or $y_3$, since properness forbids a
second $\gamma$-identity adjacent to $f_3$.  Nor is it $y_4$: in that case
the unique $\gamma$-edges at all four vertices $y_1,y_2,y_3,y_4$ would be
internal to the trunk, and, together with (7) and the missing color at
$y_0$, no trunk vertex could be incident with the assumed crossing
identity $g$.  Hence $v_1$ is outside the trunk, and the displayed
endpoint equality for $f_1$ has been obtained from its actual incidence
memberships rather than by identifying an endpoint pair with an edge.

Likewise, call the unique $\gamma$-edge at $y_4$ by the name $f_2$.
Using \noderef{Incident} at $y_4$, choose its other incident endpoint
$v_2$ and obtain
$\operatorname{ends}(f_2)=\{y_4,v_2\}$; looplessness gives
$v_2\ne y_4$.  The endpoint $v_2$ is not $y_0$, because $\gamma$ is
missing there, and it is not $y_2$ or $y_3$, because those vertices are
already incident with $f_3$.  It is not $y_1$ either: uniqueness at
$y_1$ would then give $f_2=f_1$, contradicting the just-proved fact that
the other endpoint $v_1$ of $f_1$ lies outside the trunk.  Thus $v_2$ is
also outside the trunk.  This proves the exact endpoint equality for
$f_2$ from the two incidences and looplessness.

Finally, both edge-nodes lie in the full path $P_1$.  The endpoint $y_1$
of $P_1$ is adjacent to $f_1$ in \noderef{TwoColorIncidenceGraph}, and
the edge $q$ first places $y_4$ in the component, where $y_4$ is adjacent
to $f_2$.  The same adjacencies then place $v_1,v_2$ in that component.
The vertex-containment clause of \noderef{ScheideKempePathHeightBound}
therefore places $v_1,v_2$ in $V(T)$.  These
$f_1,f_2,f_3,v_1,v_2$ satisfy every asserted endpoint, color, membership,
and outside-trunk condition.
\end{proof}
